\documentclass[../cours_1re.tex]{subfiles}
\begin{document}
\setcounter{chapter}{8}
\chapter{Calculer et justifier les dérivées}
\label{ch:09}
\begin{objectif}
Établir les règles de dérivation et les appliquer en respectant les domaines. Prérequis : nombre dérivé et continuité impliquée par la dérivabilité.
\end{objectif}
\section{Dérivée d'une somme et d'un multiple}
Si $u,v$ sont dérivables en $a$ et $\lambda$ est réel, alors
$(u+v)'(a)=u'(a)+v'(a)$ et $(\lambda u)'(a)=\lambda u'(a)$.
\begin{demonstration}
Le taux de variation d'une somme est la somme des taux ; celui de $\lambda u$ vaut $\lambda$ fois le taux de $u$. Passer à la limite donne les résultats. Une constante a un taux nul, donc une dérivée nulle.
\end{demonstration}
\section{Pourquoi la règle du produit comporte deux termes}
\begin{theoreme}
$(uv)'=u'v+uv'$ sur tout intervalle où $u$ et $v$ sont dérivables.
\end{theoreme}
\begin{demonstration}
Ajouter puis retrancher $u(a+h)v(a)$ dans le numérateur :
\[
\frac{u(a+h)v(a+h)-u(a)v(a)}h
=u(a+h)\frac{v(a+h)-v(a)}h
+v(a)\frac{u(a+h)-u(a)}h.
\]
La dérivabilité implique $u(a+h)\to u(a)$. Les deux termes tendent respectivement vers $u(a)v'(a)$ et $v(a)u'(a)$.
\end{demonstration}
\begin{attention}
La dérivée d'un produit n'est pas le produit des dérivées. Avec $u=v=x$, on aurait sinon $(x^2)'=1$, alors que la bonne valeur est $2x$.
\end{attention}
\section{Inverse et quotient}
Si $v(a)\ne0$, sa continuité assure que $v(a+h)$ reste non nul pour $h$ assez petit. Alors
\[
\frac{1/v(a+h)-1/v(a)}h
=-\frac{v(a+h)-v(a)}{h\,v(a+h)v(a)}
\longrightarrow-\frac{v'(a)}{v(a)^2}.
\]
Ainsi $(1/v)'=-v'/v^2$ ; le produit $u(1/v)$ donne
\[
\left(\frac uv\right)'=\frac{u'v-uv'}{v^2}.
\]
Le signe moins résulte de la soustraction des fractions, pas d'une convention arbitraire.
\begin{exemple}
Pour $f(x)=(x^2+1)/(x-1)$, $x\ne1$,
$f'(x)=[2x(x-1)-(x^2+1)]/(x-1)^2=(x^2-2x-1)/(x-1)^2$.
Le carré du dénominateur est positif sur le domaine ; le signe de la dérivée dépend du numérateur.
\end{exemple}
\section{Puissances entières et fonctions usuelles}
La factorisation du chapitre 2 montre $(x^n)'=nx^{n-1}$ pour $n\ge1$.
On peut aussi procéder par récurrence en utilisant le produit :
$(x^{n+1})'=nx^{n-1}x+x^n=(n+1)x^n$.
Pour $n=-m<0$, écrire $x^n=1/x^m$ et utiliser la règle de l'inverse sur $\R\setminus\{0\}$ :
$-(mx^{m-1})/x^{2m}=-mx^{-m-1}$.
Pour $n=0$, la fonction constante $1$ a pour dérivée zéro, sans devoir écrire une expression $0/x$ en zéro.
\begin{center}\begin{tabular}{lll}\toprule
Fonction & Dérivée & Domaine de dérivabilité\\\midrule
$c$ & $0$ & $\R$\\
$x$ & $1$ & $\R$\\
$x^2$ & $2x$ & $\R$\\
$x^3$ & $3x^2$ & $\R$\\
$1/x$ & $-1/x^2$ & $\R\setminus\{0\}$\\
$\sqrt x$ & $1/(2\sqrt x)$ & $]0;+\infty[$\\\bottomrule
\end{tabular}\end{center}
Chaque ligne est démontrée ici ou au chapitre 8 ; la table sert à retrouver un résultat, pas à remplacer sa justification.
\section{Composition avec une fonction affine}
Si $F(x)=f(\alpha x+\beta)$, alors $F'(x)=\alpha f'(\alpha x+\beta)$ aux points appropriés.
Pour $\alpha\ne0$, poser $k=\alpha h$ dans le taux de variation : on obtient $\alpha$ fois celui de $f$ avec accroissement $k$. Pour $\alpha=0$, $F$ est constante.
\begin{exemple}
$[(3x-1)^4]'=12(3x-1)^3$.
Sur $x>-2$, $[\sqrt{x+2}]'=1/(2\sqrt{x+2})$.
Pour $1/(2x+5)$, la dérivée est $-2/(2x+5)^2$ et $x=-5/2$ reste interdit.
\end{exemple}
\section{Complément : la règle de la chaîne}
\begin{theoreme}
Si $u$ est dérivable en $a$ et $f$ dérivable en $u(a)$, alors
$(f\circ u)'(a)=f'(u(a))u'(a)$.
\end{theoreme}
\begin{demonstration}
Écrire $u(a+h)-u(a)=h(u'(a)+\eta(h))$ avec $\eta(h)\to0$.
Écrire aussi $f(u(a)+k)-f(u(a))=k(f'(u(a))+\varepsilon(k))$, avec $\varepsilon(k)\to0$ ; définir $\varepsilon(0)=0$ rend l'identité valable même si $k=0$.
Substituer $k=u(a+h)-u(a)$ puis diviser par $h$. Le produit obtenu tend vers $u'(a)f'(u(a))$. On n'a jamais divisé par $u(a+h)-u(a)$, qui pourrait s'annuler.
\end{demonstration}
Cette démonstration générale est un approfondissement. Elle explique notamment
$(u^n)'=nu^{n-1}u'$ et $[\sqrt u]'=u'/(2\sqrt u)$ lorsque $u>0$.
\begin{finchapitre}
\section{Organiser un calcul}
Pour $F(x)=(x^2+1)\sqrt{x+2}$, le domaine est $[-2;+\infty[$ et la dérivabilité est assurée pour $x>-2$.
Le produit donne
\[
F'(x)=2x\sqrt{x+2}+\frac{x^2+1}{2\sqrt{x+2}}
=\frac{5x^2+8x+1}{2\sqrt{x+2}}.
\]
La première forme explique l'origine du calcul ; la dernière facilite l'étude du signe. Vérifier un résultat en quelques points ne démontre pas la formule, mais peut repérer une erreur.


\section{Utilité et prolongements}
Les règles permettent de transformer un modèle construit par étapes en une dérivée calculable. La règle de la chaîne suit les dépendances entre les étapes.

En optimisation et en apprentissage automatique, la différentiation automatique applique ces règles à un programme. En licence, la chaîne se généralise aux matrices jacobiennes ; en master, elle intervient dans les changements de variables et la dérivation d'applications entre espaces de fonctions. Son mécanisme reste la composition d'approximations linéaires.
\begin{retenir}
Établir le domaine, identifier l'opération principale, appliquer la règle, puis simplifier selon l'objectif. Une formule dérivée ne dispense jamais des conditions de dérivabilité.
\end{retenir}
\end{finchapitre}
\end{document}
